\documentclass[11pt,a4paper]{article}
\newcommand{\versio}{1.0}
\newcommand{\pvm}{15.10.2026}
\usepackage{fontspec}
\usepackage{polyglossia}
\setdefaultlanguage{finnish}
\setmainfont{TeX Gyre Pagella}
\setsansfont{TeX Gyre Heros}
\usepackage[a4paper,top=22mm,bottom=25mm,left=24mm,right=24mm,footskip=14mm]{geometry}
\usepackage{xcolor}
\definecolor{meri}{HTML}{2C4051}
\definecolor{leima}{HTML}{B8511F}
\usepackage{graphicx,titlesec,fancyhdr,enumitem,microtype}
\usepackage[font=small,labelfont={sf,bf,color=meri},textfont=it,labelsep=period,justification=justified,singlelinecheck=false]{caption}
\usepackage{xurl}
\usepackage[hidelinks,pdftitle={Pappila, joka ehti seistä vuoden},pdfauthor={Juri Jukola}]{hyperref}
\urlstyle{same}
\setlength{\parindent}{0pt}
\setlength{\parskip}{0.5em}
\setlength{\emergencystretch}{1.5em}
\linespread{1.03}
\titleformat{\section}{\sffamily\bfseries\large\color{meri}}{}{0pt}{}
\titlespacing*{\section}{0pt}{1.3em}{0.4em}
\newcommand{\vn}[1]{\textsuperscript{\color{leima}#1}}
\newcommand{\kuvapaikka}[1]{\fbox{\parbox[c][#1][c]{0.96\linewidth}{\centering\sffamily\small\color{leima}Kuva lisätään, kun Coden artikkelikuvat valmistuvat.}}}
\pagestyle{fancy}
\fancyhf{}
\renewcommand{\headrulewidth}{0pt}
\fancyfoot[L]{\sffamily\footnotesize\color{meri}Pappila, joka ehti seistä vuoden\enspace·\enspace versio \versio}
\fancyfoot[R]{\sffamily\footnotesize\thepage}
\fancypagestyle{eka}{\fancyhf{}\fancyfoot[R]{\sffamily\footnotesize\thepage}}

\begin{document}
\thispagestyle{eka}
{\sffamily\small\textbf{\color{leima}ARTIKKELI}\enspace{\color{leima}·}\enspace Brahea ja Pielisjärvi\enspace{\color{leima}·}\enspace versio \versio\enspace·\enspace\pvm\par}
\vspace{1.2em}
{\sffamily\bfseries\fontsize{24}{28}\selectfont\color{meri}Pappila, joka ehti seistä vuoden\par}
\vspace{0.6em}
{\sffamily\large\color{meri}Missä Mårten Liukkosen talo oli?\par}
\vspace{0.5em}
{\sffamily\small Juri Jukola\par}
\vspace{1.0em}
{\color{leima}\rule{\linewidth}{0.8pt}}\par
\vspace{0.3em}
{\itshape
Kesällä 1655 talonpoika Mårten Liukkonen luovutti talonsa, jotta Brahean kirkkoherralla olisi pappila kirkon vieressä. Seuraavana kesänä vihollinen poltti sen. Talon paikkaa ei ole kukaan etsinyt. Asiakirjat, vanhat kartat ja laserkeilaus rajaavat sen nyt kahteen alueeseen Varasjoen suulla, joen kummallekin rannalle.\par}
\vspace{0.4em}
{\color{meri}\rule{\linewidth}{0.4pt}}\par

\section{Kaksi yötä metsässä}

Brahean talvikäräjillä tammikuussa 1680 kirkkoherra Lars Laurentii pyysi todistusta siitä, miten hän oli hoitanut virkaansa. Hän oli ollut seurakunnassa kuusi vuotta kappalaisena ja sen jälkeen 20 vuotta kirkkoherrana. Lautakunta ja koko rahvas todistivat, että hän oli elänyt sovussa sekä esimiehensä että seurakuntalaistensa kanssa ja antanut köyhille anteeksi heidän maksunsa. Sitten kirjuri merkitsi pöytäkirjaan jotain, mitä käräjillä harvoin näki. Kirkkoherra kertoi asiansa \emph{medh gråtande tåhrar} (\emph{kyynelsilmin}).\vn{1}

Lars Laurentii muisteli vuotta 1656. Kun venäläiset hyökkäsivät Pielisjärvelle, hänen nyt jo edesmennyt vaimonsa joutui kahden pienen lapsen kanssa olemaan kaksi yötä metsässä \emph{blott och baar} (\emph{paljaana ja puutteessa}). Lars itse seisoi talonpoikien kanssa vihollista vastaan, ja vihollinen \emph{Kyrckian, Prästegården och några andra bonde gårdar afbrände} (\emph{poltti kirkon, pappilan ja muutamia talonpoikaistaloja}). Lopulta vihollinen ajettiin pakoon hänen neuvostaan. Kirkkoherran todisteena oli konsistorin kirje Viipurista maaliskuulta 1657. Pappila, joka paloi 1656, oli ollut pappilana vain vuoden. Se oli Mårten Liukkosen talo. Tässä kirjoituksessa yritän selvittää, missä se seisoi.

\section{Kirkkoherra ei halunnut soutaa}

Brahean kaupunki perustettiin 1653 Lieksanjoen suulle, ja samalla määrättiin rakennettavaksi kirkko, käräjätupa ja pappila. Pappilan paikasta käytiin käräjillä kahden vuoden ajan keskustelua. Talvella 1654 sopivimpana pidettiin Michofara Larisen taloa, joka oli ollut pappila jo ennen, mutta kreikanuskoisen papin pappila. Larisen talo oli Lieksanjoen toisella puolella.\vn{2}

Kesäkäräjillä 1654 kirkkoherra Anders pyysi toista taloa. Hän halusi Mårten Liukkosen talon, \emph{alldenstund det blifwer närmare kyrkian beläget} (\emph{koska se on lähempänä kirkkoa}), eikä Larisen talolta keväisin päässyt joen yli sille puolelle, jolle kirkkoa rakennettiin.\vn{2} Kirkkoherra ei halunnut jäädä keväisin joen taakse. Hänen pyyntönsä on ensimmäinen vihje talon paikasta: se oli kirkon lähellä ja samalla puolella Lieksanjokea.

Kesäkuun 20. päivänä 1655 käräjät päättivät, että \emph{Liuckanens hemman förordnas till Prästebool uthi Brahea} (\emph{Liukkosen talo määrätään Brahean pappilaksi}). Liukkonen luovutti talonsa vapaaehtoisesti ja sai tilalle Larisen talon. Porvari Matz Caupinen oli jo ehtinyt ostaa Larisen talon, joten pitäjä maksoi hänelle rahat takaisin. Samoilla käräjillä rahvasta muistutettiin kirkon, pappilan ja käräjätuvan rakentamisesta Brahean kaupunkiin.\vn{3} Pöytäkirja ei kerro, rakennettiinko Liukkosen talolle uusia rakennuksia vai otettiinko se käyttöön sellaisenaan.

\section{Vanha kolmen talon paikka}

Mårten Liukkonen ei ollut tulokas. Kesäkäräjillä 1653 hän valitti, että käskynhaltija oli talvella perinyt häneltä viisi hopeataaleria. Liukkonen oli viljellyt oman peltonsa vieressä autiotalon peltoa, ja entinen pitäjänkirjuri muisti, että Liukkonen oli ollut samasta pellosta velkaa jo 19 vuotta aiemmin. Liukkonen oli siis viljellyt maata kirkon lähellä ainakin 1630-luvulta.\vn{4}

Lautakunta selvitti samalla, mistä paikasta oli kyse: \emph{der be:te Liuckanen och Maxima Sawanpoica nu boo, hafwa tillförende trenne bönder bodt} (\emph{siellä, missä Liukkonen ja Maxima Savanpoika nyt asuvat, on ennen asunut kolme talonpoikaa}). Liukkonen ja Maximan isä olivat ottaneet kolmen talon maat ja jakaneet ne kahtia, ja Liukkosella oli suurin osa. Käräjät antoivat kolmannen talon Staffan Lodwoiselle, ja maat jakoivat Nils Sipolan, Oluf Partainen ja Lauri Sabattinen.\vn{4} Kahden vuoden päästä samassa riidassa Maximan poika on kirjattu nimellä \emph{Sawa Maximanpoica Sabattinen}.\vn{5} Liukkosen naapurit olivat Sapattisia. Tämä on toinen vihje. Liukkosen talo oli Lieksan kylän vanhaa ydintä, kolmen talon paikalla, jonka se jakoi Sapattisten kanssa.

\section{Talo, joka ei kadonnut}

Pappilan polton jälkeen kirkkoherran virkatalo siirtyi viimeistään 1659 Lieksanjoen toiselle puolelle Laristen entisiin taloihin, kun kreikanuskoiset Larit olivat lähteneet sodan aikana. Siitä tuli Pielisjärven iso pappila, joka pysyi paikallaan vuosisatoja.\vn{6} Liukkosten vaiheet näkyvät Kajaanin vapaaherrakunnan maakirjoista, joihin kirjattiin vuosittain jokainen veroa maksava talo.

Vuonna 1655 Lieksan kylän luettelossa on yksi Liukkonen. Vuonna 1656, pappilan palovuonna, luetteloon ilmestyy uusi rivi: \emph{Michell Liukkainen}, heti Maxima Savanpojan jälkeen. Sama järjestys toistuu vielä 1670-luvulla. Vuosina 1670, 1675 ja 1677 \emph{Michali Liukkonen} on joka kerta heti Sava Maximan ja Hoatta Savanpojan jälkeen. Vuonna 1677 kirjuri merkitsi hänen kohdalleen verohuojennuksen, koska \emph{frost fortagit hans åhrswäxt} (\emph{halla vei hänen satonsa}).\vn{7} Kun Lars Laurentii itki käräjillä 1680, lautakunnassa istui Hoatta Savanpoika, Liukkosten naapuri.\vn{1}

Liukkoset eivät siis muuttaneet pois kaupungin ytimestä. Mihail eli Michali Liukkonen, todennäköisesti Mårtenin poika, isännöi taloa Sapattisten vieressä yli 20 vuotta palon jälkeen. Kantatilaluettelossa, joka on koottu samoista maakirjoista vuosisatoja myöhemmin, isäntä on nimeltään Mikitta Liukkonen, ja talon kerrotaan siirretyn \emph{kirkon paikalta} Timitranniemeen 1650-luvulla.\vn{8} Maakirjoissa nimi on kuitenkin joka vuosi Michell tai Michali, ja talo pysyy kaupungin ryhmässä. Luettelo yhdistää Liukkosten talon ja toisen talon, jonka isännät Hemming Härkönen ja Sigfrid Heinonen ovat maakirjoissa omana rivinään. Siirto Timitranniemeen koskee todennäköisesti tätä toista taloa tai myöhempää aikaa. Ilmaus \emph{kirkon paikalta} on silti kolmas vihje: talo oli kirkon tontilla tai sen vieressä.

\section{Sapattisten naapurina}

Sapattisten talot liitettiin 1685 Brahen hoviin, kun kaupunki oli lakkautettu ja sen maat koottiin vuokrattavaksi kartanoksi, Hovilaksi. Vuosien 1838–1840 isojakokartalla Hovilan maat ovat vanhan kirkkomaan ja käräjäpaikan kohdalla Varasjoen molemmin puolin, ja käräjäpaikka, \emph{Tingsplan}, on merkitty pohjoisrannan entisen kirkkotontin viereen.\vn{9} Liukkosen talo oli Sapattisten vieressä, joten sekin oli tällä alueella.

Vuosien 1818–1819 verollepanokartta näyttää Varasjoen suun kymmenen vuotta ennen kuin kirkko paloi 1828. Kirkon ja hautausmaan tontti on Varasjoen pohjoisrannalla, kellotapulin tontti etelärannalla, ja niiden välillä on silta.\vn{10} Muu maa tonttien ympärillä on numeroituja palstoja. Kun maa on 1700-luvulla kynnetty pelloksi, kyntö on tasoittanut talonpaikoista kuopat ja matalat kummut. Säilyä voi vain kiviraunio, jota auran kanssa ei kannattanut kiertää.

\begin{figure}[t]
\centering
\includegraphics[width=\linewidth]{kuva_A_sijaintikartta.png}
\caption{Liukkosen talon mahdolliset alueet: A Varasjoen etelärannalla tapulin tontin ja Kirkkosaari~1:n muinaisjäännösalueen ympärillä, B pohjoisrannalla, jossa olivat myöhemmät kirkot ja vuosien 1838–1840 isojakokartan käräjäpaikka. Tarkka paikka on tuntematon. Ympyrä merkitsee kumpua, joka erottuu tarkemmassa laserkeilausaineistossa mutta ei tämän kuvan pohjassa. Tonttimaat on sijoitettu vuosien 1818–1819 verollepanokartalta noin 30 metrin tarkkuudella. Sisältää Maanmittauslaitoksen Laserkeilausaineisto 0,5 p -aineistoa ja maastotietokannan aineistoa (CC BY 4.0). Kirkkosaari~1: muinaisjäännösrekisteri, Museovirasto. Tulkinta: J.~Jukola 2026.}
\end{figure}

\section{Kaksi mahdollista aluetta}

Vihjeet rajaavat talon noin 150 metrin säteelle Varasjoen suusta. Sen sisällä on kaksi aluetta, joen kummallakin rannalla, ja kummallekin on omat perusteensa.

\textbf{A. Varasjoen eteläranta tapulin tontin ympärillä.} Paikka oli kirkon lähellä, vain kapean joen ja sillan takana, ja samalla puolella Lieksanjokea. Kirkkosaari~1 -niminen kohde on muinaisjäännösrekisterissä historiallinen asuinpaikka, ja sen koekuopasta löytyi 2014 tiiltä ja hiiltä.\vn{11} Tiiltä käytettiin 1600-luvulla säätyläisten taloissa ja pappiloissa, ei talonpoikien tuvissa. Laserkeilauksessa alueella on yksi selvä kumpu, noin viisi metriä leveä ja metrin korkea, 26 metrin päässä tapulin tontista. Se on pyöreä ja symmetrinen, eikä sen ympärillä ole vallia tai kuoppaa. Se on kaikista alueen kohteista selvin, mutta laser ei kerro, onko se palaneen talon uunin raunio vai pellolta kerättyjen kivien kasa. Ympäriltä ei löydy muita talonpaikan merkkejä, mutta kynnetyllä pellolla sellaisia ei odotakaan. Vuosien 1818–1819 verollepanokartalla kohta ei ole tonttimaata, eikä kartan väristä voi päätellä, oliko se peltoa vai hakaa. Vuoden 1956 ilmakuvassa paikka on avointa sarkapeltoa, mutta kumpua ei erota kuvasta. Se on voinut olla pellon keskellä kiertämättä jääneenä kasana, tai se on syntynyt vasta myöhemmin.

\textbf{B. Varasjoen pohjoisranta: kirkon paikka ja käräjäpaikka.} Tälle alueelle osoittaa kaksi erillistä vihjettä. Kantatilaluettelon mukaan Liukkosen talo siirrettiin kirkon paikalta. Jos Brahen kirkko rakennettiin 1667 sille kohdalle, jolla Liukkosen talo oli seissyt, talo oli tällä matalalla kohoumalla, jota asiakirjat kutsuvat saareksi. Toinen vihje tulee naapureista: Sapattisten talot siirtyivät Hovilaan, ja Hovilan käräjäpaikka oli entisen kirkkotontin vieressä. Laserkeilaus ei näytä kielekkeellä talonpohjia, mutta se ei näytä siellä kirkkojenkaan pohjia, joten maa kätkee täällä kaiken alle 20 sentin korkuisen.

Itse pidän etelärantaa todennäköisimpänä, mutta se on päätelmä eikä asiakirjatieto. Varmaa on, että talo oli Brahean ytimessä kirkon lähellä ja että sen paikka on edelleen kynnetyn maan alla.

\begin{figure}[t]
\centering
\includegraphics[width=\linewidth]{kuva_B_verokartta_1818_2x.png}
\caption{Varasjoen suu vuosien 1818–1819 verollepanokartalla. Ylemmän sillan molemmin puolin ovat kirkon ja tapulin punaiseksi väritetyt tonttimaat. Ympyrät osoittavat laserkeilauksessa havaitun kummun ja nykyisen kirkon kohdan noin 30 metrin tarkkuudella. Kummun kohdalla on vuonna 1818 palstan 10 numero eikä tonttimaata. Kansallisarkisto, Maanmittaushallituksen verollepanokartat, 125:1 Lieksa 4:2a–l, lehti 8; merkinnät J.~Jukola.}
\end{figure}

\section{Miten paikka löytyy}

Kaksi asiaa voi ratkaista kysymyksen. Isojaon jakokirja vuosilta 1840–1847 on tilattu Kansallisarkistosta. Siitä näkee, millä palstoilla Hovilan maat olivat ja mikä palsta oli Sapattisten talon. Liukkosten talo oli sen vieressä kaksikymmentä vuotta. Toinen tie kulkee maaston kautta. Yksi käynti kummulla kertoo, onko se pellon kiviä vai palaneita uunikiviä, savilaastia ja tiiltä. Kumpaakaan ei voi tehdä kotoa käsin.

Pappila seisoi vain vuoden, ja sen viimeisestä yöstä on jäljellä yksi lause kirkkoherran todistuksessa neljännesvuosisataa myöhemmin. Talon isäntä ja hänen poikansa jäivät kuitenkin samalle paikalle, maksoivat veronsa ja menettivät satonsa hallalle kuten muutkin. Heidän talonsa paikka on Lieksan ensimmäisen pappilan paikka, ja se on löydettävissä.

\section{Lähteet}

{\footnotesize
\begin{enumerate}[label=\arabic*,leftmargin=1.6em,itemsep=0.02em,topsep=0.1em]
\item Kansallisarkisto (KA), Kajaanin vapaaherrakunnan renovoidut tuomiokirjat, KO a:2, talvikäräjät Braheassa 3.\,ja 5.1.1680, kuvat 429–430.
\item KA, Kajaanin vapaaherrakunnan renovoidut tuomiokirjat, KO a:1, talvi- ja kesäkäräjät 1654, kuvat 197–199.
\item KA, KO a:1, käräjät 20.6.1655, kuva 258.
\item KA, KO a:1, kesäkäräjät 1653, kuva 134.
\item KA, KO a:1, käräjät 1655, kuva 230.
\item Kantatilaluettelo, Isopappila; KA, Kajaanin vapaaherrakunnan maakirja 1657.
\item KA, Pohjanmaan läänin tilejä, Kajaanin vapaaherrakunnan maa- ja tilikirjat 1655 (s.~27), 1656 (s.~30), läänintili- ja maakirjat 1670 (s.~378), 1675 (s.~274) ja 1677 (s.~102).
\item Kantatilaluettelo, N:o 69 Dimitra.
\item Kantatilaluettelo, N:o 77 Hovila; ote A.\,J. von Fieandtin isojakokartasta 1838–1840 teoksessa Poutiainen, Sepänmaa \& Jussila 2019: \emph{Lieksan keskustaajaman osayleiskaavan arkeologinen inventointi}, s.~14.
\item KA, Maanmittaushallituksen verollepanokartat, 125:1 Lieksa 4:2a–l, C.\,G. Backman 1818–1819, lehti 8.
\item Museovirasto, muinaisjäännösrekisteri: Lieksa, Kirkkosaari 1; inventointi V.~Laakso, J.~Ruohonen ja S.~Kärkkäinen 2014.
\end{enumerate}

\vspace{0.4em}
{\color{leima}\rule{\linewidth}{0.4pt}}\par
{\sffamily Juri Jukola · jujukol@gmail.com · linkedin.com/in/jurijukola}\par
Artikkeli on vapaasti luettavissa ja jaettavissa lähde mainiten. Korjaukset ja lisätiedot ovat tervetulleita.\par
}
\end{document}
